\chapter{Discussion}
\label{chap:discussion}

\section{Interpretation of Error Type Findings}
\label{sec:error-type-interpretation}

The most noticeable pattern in the error-type data is not how the models
differ, but what they have in common. Combinatorial Counting Error is
the single most common error type for every model individually
(Section~\ref{sec:rq1}): 15 of Gemini's 45 incorrect responses
(33.3\%), 9 of GPT-OSS-120B's 36 incorrect responses (25.0\%), and 20 of
Mistral Large's 45 incorrect responses (44.4\%). In each case this is
clearly ahead of that model's own next most common error type
(Incomplete Reasoning for Gemini and Mistral Large; Logical Reasoning
Error for GPT-OSS-120B). This shows that counting problems -- such as
over- or under-counting, or confusing permutations with combinations --
are a shared structural weakness for Gemini, GPT-OSS-120B, and Mistral
Large. It is not a problem unique to any one model or provider.

This shared weakness should be considered alongside the Chi-square test
result in Section~\ref{sec:rq1} ($\chi^2 = 28.35$, $p = 0.101$), which
does not show statistically significant evidence that the overall
distribution of error types differs across models. Any model-specific
patterns discussed below should be seen as descriptive observations, not
statistically confirmed differences, especially given the small expected
cell counts noted in that section.

With that in mind, one descriptive pattern stands out. On the Hard-tier
common-wrong-question set (Section~\ref{sec:hard-tier-freq}),
GPT-OSS-120B's error profile under CoT is different from the other two
models. Its most common error is Logical Reasoning Error, not
Combinatorial Counting Error. Under ToT, this difference goes away, and
all three models show Combinatorial Counting Error as their most
frequent failure on the same question set. One possible explanation is
that ToT's branch-and-select structure filters out some types of errors
better than others. For example, a logical error in one branch can be
avoided by choosing another branch. Still, a miscounting error that
happens in all reasoning approaches cannot be fixed just by generating
and selecting between multiple attempts. This suggests that
counting-type errors are a deeper, more structural limitation, while
logical-reasoning mistakes are more likely to be caught through
reflection or resampling. However, this interpretation is speculative
and is not directly tested by this study.

\section{Interpretation of Step Count and Correctness Findings}

The step-count pattern described in Section~\ref{sec:step-count} is one
of the few areas where Mistral Large acts differently from the other two
models, and this difference may be meaningful. As explained in
Section~\ref{sec:models}, Gemini and GPT-OSS-120B both have an internal
reasoning-effort control set to its lowest level and kept constant
during this study. Mistral Large does not have this control and is used
as a standard instruct model. As a result, Gemini's and GPT-OSS-120B's
reasoning length is produced under a constant internal budget setting,
while Mistral Large's reasoning length is free to vary as much as the
model chooses on a given question.

This difference helps explain why only Mistral Large shows a gap in step
counts between correct and incorrect Hard-tier answers (mean 12.59 vs.
16.12; median 7.0 vs. 12.0), while Gemini and GPT-OSS-120B do not. If a
model's internal reasoning-effort setting is pinned low and held
constant, its reasoning length may have little additional room to expand
specifically on the questions it struggles with, so correct and
incorrect answers end up costing a similar number of steps. Mistral
Large, without this constraint, is free to generate a longer, more
exploratory chain of steps precisely on the questions it finds
difficult, and this extra effort does not appear to improve its accuracy
-- if anything, it often takes more steps on the questions it gets wrong
than on the ones it gets right.

Another possible explanation is that some Hard-tier questions are more
difficult than others. These harder questions might require more steps
from any model and are more likely to be answered incorrectly. If this
were the main reason, a similar step-count gap would be expected for
Gemini and GPT-OSS-120B, since all three models faced the same 40
Hard-tier questions. Since this gap does not appear for the other two
models, the pure-difficulty explanation seems less convincing as the
only cause. Still, it may play a role in how Mistral Large handles
difficult questions. Like RQ2 itself, this is an interpretation of the
data, not a proven cause.

\section{Interpretation of CoT vs.\ ToT Findings}

Section~\ref{sec:overall-accuracy} (CoT vs. ToT) shows that ToT improves
accuracy significantly for two models, GPT-OSS-120B and Mistral Large,
but not for Gemini. Section~\ref{sec:tot-design} explains that ToT
requires about four times as many generation calls as CoT, and
Section~\ref{sec:execution-time} shows the practical result: mean execution
time increases by 2.29x for Gemini, 4.62x for GPT-OSS-120B, and 2.99x
for Mistral Large. Whether this extra cost is justified depends on the
model, and the answer is different for each one.

For GPT-OSS-120B and Mistral Large, ToT is a strong option. Both models
show a clear accuracy improvement in the Hard tier, where improvement is
needed (Section~\ref{sec:overall-accuracy}): GPT-OSS-120B gains 20.0
percentage points (45.0\% to 65.0\%) and Mistral Large gains 17.5
percentage points (37.5\% to 55.0\%). GPT-OSS-120B has the biggest
relative cost increase (4.62x, close to the expected 4x). Still, it is
also the model with a significant McNemar result, which suggests that
the extra computation leads to a real benefit.

For Gemini, there is no case to make at all. Its Hard-tier accuracy is
identical under both strategies (50.0\% for CoT and ToT), and its
Medium-tier accuracy actually falls slightly under ToT (100.0\% to
97.5\%). With only 6 discordant pairs, split evenly 3 to 3
(Section~\ref{sec:overall-accuracy}), McNemar's test gives $p = 1.0$ --
about as clean a null result as this kind of test can produce. Gemini
still pays a 2.29x execution-time cost for ToT, without any accuracy
benefit to offset it, demonstrated or otherwise. It is worth being
precise about what this non-significant result does and does not show:
it does not prove that ToT has no effect on Gemini in principle, only
that in this study's data, there is no detectable one, and what
difference does exist points slightly the other way. A larger sample of
discordant pairs might yet disclose a significant difference in either
direction, but the current data give no reason to expect one favouring
ToT.

Overall, these results do not give a single answer to whether ToT is
worth its cost. For two of the three models, ToT gives a clear and
meaningful accuracy boost that likely justifies its 3--4.6x cost. For
the third model, the cost is paid without any accuracy benefit at all,
proven or otherwise. For any new model, this question should be tested
directly, just as it was here.

\section{Threats to Validity}
\label{sec:threats-validity}

Section~\ref{sec:limitations} presents the main methodological
limitations of this study. This section explains how each limitation
affects the conclusions in Chapter~\ref{chap:results}, instead of
repeating the list.

The \textbf{reasoning-effort confound} is most important for RQ3.
GPT-OSS-120B's reasoning-effort setting is set to low, not zero, and the
model is designed to reason, so some internal reasoning may still happen
even without explicit CoT/ToT prompts. If GPT-OSS-120B does more internal
reasoning at low effort than Gemini does at its ``minimal'' thinking
level, this difference -- not just the prompting structure -- could help
explain why GPT-OSS-120B shows a significant ToT improvement while
Gemini does not (Section~\ref{sec:overall-accuracy}). This study cannot
separate these two explanations, because the internal reasoning budget
was only kept constant within each model's own CoT/ToT comparison, not
matched across models.

\textbf{Data contamination risk} threatens a more basic assumption: that
accuracy differences show reasoning ability, not just memorised training
data. This matters for the near-perfect Easy/Medium accuracy in
Section~\ref{sec:overall-accuracy} (95--100\% for every model and
prompting strategy). Another possible reason for this ceiling effect is
that Hendrycks MATH, a common public benchmark, is likely included in
the pretraining data for all three models. High accuracy may therefore
reflect recall instead of real problem-solving. This study cannot tell
these explanations apart, which is why the Hard tier, based on the
less-republished AIME, is more important for analysis in this
dissertation.

\textbf{Judge conflict of interest} threatens RQ1 specifically: Gemini
is both a model under evaluation and the model classifying every
incorrect response's error type. The planned mitigation, an independent
sample re-judged by a neutral model with agreement measured by Cohen's
kappa, was not completed (Section~\ref{sec:judging}), so this bias is
neither confirmed nor ruled out. Some RQ1 findings are more exposed to
this risk than others: a claim that Gemini's own errors are classified
differently from the other two models' errors would be directly
vulnerable to judge leniency, whereas the shared-weakness finding
discussed in Section~\ref{sec:error-type-interpretation} (Combinatorial Counting Error
dominant across all three models, including the two Gemini has no
evaluative stake in) is comparatively less exposed, since a bias
favourable to Gemini specifically could not manufacture a pattern shared
by all three models' outputs.

The \textbf{fixed Tree-of-Thought branch count} of three limits how much
the RQ3 conclusions can be generalised. The finding that ToT
significantly improves accuracy for two of the three models only applies
to this study's three-branch setup, not to tree-search-based prompting
in general. Using a different branch count could lead to different
results for any model.

Finally, \textbf{small subgroup sample sizes} limit how much confidence
can be placed in the more detailed breakdowns in this dissertation. The
Hard tier's category counts range from 8 (Probability) to 16 (Algebra
and Combinatorics), and RQ1's Chi-square test was flagged for a low
minimum expected cell count of 0.286. Claims at the category level or
for each model in each category should be seen as descriptive, not
statistically confirmed, as already noted in Sections~\ref{sec:rq1}
and~\ref{sec:limitations}.

\section{Comparison with Prior Work}

Sun et al.~\cite{sun2025errorclassification} and Zhang and
Graf~\cite{zhang2025mathematical} both classify LLM mathematical
reasoning errors within individual models. By comparing the same error
taxonomy across three models, this dissertation finds something those
single-model studies could not: the largest error category, Combinatorial
Counting Error, appears in all three models, not just one. The
Chi-square test also shows no statistically significant difference in
error-type distribution between models (Section~\ref{sec:rq1}). While
Sun et al. and Zhang and Graf show that error type and location are
measurable and informative, this dissertation adds that at least one
property, the dominant error type, may be common across models rather
than unique to each.

Section~\ref{sec:tot-scope} explains that the Tree-of-Thought method
tested here is a smaller-scale version of what Yao et
al.~\cite{yao2023tree} describe. Instead of a deep, iterative search,
this study uses three independently generated candidates and one
selection step. The RQ3 result -- that this simpler approach improves
accuracy for two models but not Gemini -- applies only to this limited
version, not to all tree-search-based prompting. It is still unclear
whether the full iterative search from Yao et al. would help Gemini or
further benefit the other models. Wang et al.'s
Self-Consistency~\cite{wang2023selfconsistency}, which combines
independently sampled CoT
paths by majority vote, gives a similar warning: more reasoning does not
always mean more reliable answers if the same error appears in each
sample. The RQ2 finding here, that Mistral Large's longer,
less-constrained reasoning on the Hard tier does not lead to higher
accuracy, supports this warning, even though the two studies use
different methods.

Finally, Hao et al.~\cite{hao2024autorace} point out the risk of
self-evaluation bias when an LLM judges its own or similar models'
outputs. This dissertation faces that risk directly, since Gemini is
both being evaluated and is the main error-classification judge.
Section~\ref{sec:threats-validity} looks at how this risk affects the
RQ1 results in more detail. The key point here is that this dissertation
tried to reduce the risk in the way Hao et al. suggest -- by using an
independent, neutral judge for validation -- though this was not done at
the planned scale (Section~\ref{sec:judging}).

\section{Summary of Findings}
\label{sec:summary-findings}

RQ1 looked at whether the types of mathematical reasoning errors varied
between the three models. The most common error, Combinatorial Counting
Error, appeared in all three models -- Gemini, GPT-OSS-120B, and Mistral
Large -- so it was not unique to any one of them. A Chi-square test
showed no significant difference in how error types were distributed
across the models. There was a shift in the Hard-tier
common-wrong-question set: for GPT-OSS-120B, Logical Reasoning Error was
most common under CoT, but this changed under ToT, where all three
models showed the same main error type.

RQ2 explored whether the number of reasoning steps a model takes is
related to whether its answer is correct. For Gemini and GPT-OSS-120B,
there was no clear link on the Hard tier -- both correct and incorrect
answers used about the same number of steps. For Mistral Large, however,
incorrect answers on the Hard tier used more steps than correct ones
(mean 16.12 vs. 12.59; median 12.0 vs. 7.0). This may be because Mistral
Large does not have a way to limit its reasoning length, unlike the
other two models.

RQ3 looked at whether Tree-of-Thought (ToT) prompting leads to better
answers than Chain-of-Thought (CoT) for each model. McNemar's Exact Test
showed that ToT gave a significant improvement for GPT-OSS-120B
($p = 0.021$) and Mistral Large ($p = 0.016$), but not for Gemini
($p = 1.0$). The improvement was mainly seen in the Hard tier, where
there was more room to get better results. For Easy and Medium
questions, accuracy was already very high for all models. Therefore, ToT
is not equally helpful for every model, and its much longer execution
time is only worth it when there is a clear accuracy gain.

\section{Future Work}

Several directions for future work follow directly from the limitations
discussed in Section~\ref{sec:threats-validity}.

The Tree-of-Thought implementation in this study used a fixed branch
count of three, chosen for scope rather than tested against alternatives
(Section~\ref{sec:threats-validity}). A natural extension is to test a
range of branch counts to establish whether the significant accuracy
gains seen for GPT-OSS-120B and Mistral Large continue to grow with more
branches, plateau, or whether a larger branch count could close the gap
for Gemini, where no significant gain was found at three branches.

The planned neutral-judge validation was not completed at the intended
scale of 120 responses due to the neutral judge provider's daily token
quota (Section~\ref{sec:judging}). Completing this validation with more
time or a higher-quota provider would directly address the judge
conflict of interest discussed in Section~\ref{sec:threats-validity},
and would clarify whether Gemini's dual role as both evaluated model and
judge introduces a measurable bias in the RQ1 error-type findings.

Similarly, the planned manual grading validation of the automatic answer
evaluator against human judgment, covering a sample of real model
outputs rather than only hand-picked test cases, was not carried out
(Section~\ref{sec:limitations}). This would strengthen confidence in the
answer-correctness labels underlying every result in this dissertation.

This study evaluated three free-tier models, none of which are dedicated
frontier reasoning models with reasoning always on. Extending this
study's design to such models would directly test the reasoning-effort
confound discussed in Section~\ref{sec:threats-validity}, since a model
without a low-effort setting to pin would remove one source of
uncertainty in interpreting the RQ3 results.

The Hard tier's AIME-only category counts are unbalanced (Algebra 16,
Combinatorics 16, Number Theory 11, Probability 8), limiting the
statistical power of category-level comparisons
(Section~\ref{sec:threats-validity}). A larger, rebalanced Hard-tier
sample, curated and categorised beyond what was available for this
project, would allow category-level claims to be tested with confidence
rather than treated as descriptive.

This dissertation's RQ2 was limited to a descriptive analysis in the
main text (Section~\ref{sec:step-count}), though an exploratory
Mann-Whitney U test and logistic regression were run as a supplementary
check, broken down by model, prompt, and difficulty tier. The results
were mixed and test-dependent: for example, Gemini's Hard-tier CoT gap
reached significance under the Mann-Whitney test ($p = 0.008$) but not
under logistic regression ($p = 0.219$) on the same data, and no other
model/prompt cell reached significance under either test. Because the
two tests disagreed on the one notable case, and most cells were
underpowered at this sample size, these results were not treated as a
confirmed finding in the main text. A larger-scale replication, with
enough data to give consistent results across test choices, remains a
natural next step.

The Tree-of-Thought design used here creates three candidates
independently and then picks one in a single round. This is different
from the deeper, step-by-step evaluation and pruning described by Yao et
al.~\cite{yao2023tree} (Section~\ref{sec:tot-scope}). Implementing this
more thorough iterative search, and also comparing it to Wang et al.'s
Self-Consistency~\cite{wang2023selfconsistency} (which uses independent
sampling and majority voting, but is simpler to engineer than
branch-and-select ToT), could help show whether ToT's advantage comes
from evaluating and choosing between branches, or just from making
multiple attempts at a problem.

Combinatorial Counting Error is the biggest error category found in all
three models (Section~\ref{sec:rq1}). A follow-up study could create a
smaller, targeted diagnostic set to focus on specific counting
sub-skills, for example choosing between permutations and combinations,
using inclusion-exclusion, and spotting over- or under-counting from
double-counted cases. This would help identify which sub-skill causes
this shared weakness, instead of treating ``combinatorial counting'' as
one broad category.

Finally, the results of this study depend on specific model versions at
a certain point in time (Section~\ref{sec:models}). Model names like
\texttt{gemini-flash-lite-latest} and \texttt{mistral-large-latest} can
point to different models if the study is repeated later. Repeating this
research with newer model versions would show whether the patterns found
here -- especially the shared Combinatorial Counting Error weakness and
the model-dependent ToT benefit -- continue as models improve or are
unique to this generation.

